\documentclass[UTF8]{ctexart}   % UTF8 编码，使用 ctex 内置字体
\usepackage{geometry}           % 设置页面边距
\usepackage{graphicx}            % 插入图片
\usepackage{amsmath, amssymb}    % 数学公式
\usepackage{hyperref}            % 超链接

\geometry{a4paper, left=2.5cm, right=2.5cm, top=2.5cm, bottom=2.5cm}

\newcommand{\ptl}[2]{\frac{\partial #1}{\partial #2}}

\begin{document}

\section{RNN网络(深度平衡)}
    \subsection{理论推导}
        设RNN网络的迭代公式为：
        \begin{equation}
            h_{t+1} = \underset{h}{argmin}~E(h, h_t, x_t;\theta)
        \end{equation}
        其中：
        \begin{itemize}
            \item $h_t$为t时刻的隐变量
            \item E为系统的能量函数
            \item $x_t$为系统t时刻的输入
            \item $\theta$为系统的可训练参数
        \end{itemize}

        记RNN网络的读出和误差为：
        \begin{align}
            &y_t = \phi_{out}(h_t)\\
            &Loss = L_t(y_t, \hat{y}_t)
        \end{align}

        记$E_t := E(h_t, h_{t-1}, x_t;\theta)$,则有:
        \begin{equation}
            g_t(h_t, h_{t-1};\theta) := \ptl{E_t}{h_t} = \mathbf{0}。
        \end{equation}

        上式取全微分可得：
        \begin{equation}
            0 = \ptl{^2 E_t}{h_t^2} dh_t + \ptl{^2E_t}{h_t \partial h_{t-1}} dh_{t-1} + \ptl{^2 E_t}{h_t \partial \theta} d\theta
        \end{equation}

        我们定义：
        \begin{align}
            A_t &:= -\left(\ptl{^2 E_t}{h_t^2}\right)^{-1}\\
            B_t &:= \ptl{^2 E_t}{h_t \partial h_{t-1}}\\
            d_\theta g_t &:= \ptl{^2 E_t}{h_t \partial \theta} d\theta\\
        \end{align}

        则有：
        \begin{equation}
            dh_t = A_t B_t dh_{t-1} + A_t d_\theta g_t
        \end{equation}

        损失函数的全微分为：
        \begin{align}
            dL_t &= \ptl{L_t}{y_t} dy_t = \ptl{L_t}{y_t} \phi_{out}'(h_t) dh_t \\
            &:= v_t^T dh_t \\
            &= v_t^T A_t d_\theta g_t + v_t^T A_t B_t dh_{t-1}
        \end{align}

        对于任意给定的时刻$\tau$，我们定义如下伴随向量：
        \begin{align}
            p_\tau^T &:= v_\tau^T A_\tau \\
            p_{t-1}^T &:=p_\tau^T B_t A_{t-1}
        \end{align}
        则有：
        \begin{equation}
            \begin{aligned}
                dL_\tau &= p_\tau^T d_\theta g_\tau + p_\tau^T B_\tau dh_{\tau-1} \\
                &= p_\tau^T d_\theta g_\tau + p_\tau^T B_\tau A_{\tau-1} d_\theta g_{\tau-1} + p_\tau^T B_\tau A_{\tau-1} B_{\tau-1} dh_{\tau-2} \\
                &= p_\tau^T d_\theta g_\tau + p_{\tau-1}^T d_\theta g_{\tau-1} + p_{\tau-1}^T B_{\tau-1} dh_{\tau-2} \\
                &\dots \\
                &= \sum_{t=1}^{\tau} p_t^T d_\theta g_{t}
            \end{aligned}
        \end{equation}
        从而损失函数关于$\theta$的梯度为：
        \begin{equation}
            \label{eq:1}
            \frac{\partial L_\tau}{\partial \theta} = \sum_{t=1}^{\tau} \left( \frac{\partial g_t}{\partial \theta} \right)^T p_t = \sum_{t=1}^{\tau} \left(\ptl{^2 E_t}{h_t \partial \theta}\right)^T p_t
        \end{equation}

    \subsection{算法概述}
        \paragraph{向前传播}
        利用优化器求解$h_t = argmin~E(h_t, h_{t-1}, x_t; \theta)$，并保存$h_t$
        \paragraph{损失计算}
        根据当前输出计算损失$Loss = L_\tau(y_\tau, \bar{y}_\tau)$， 然后计算列向量$v_\tau$：
        \begin{equation}
            v_\tau = \left( \frac{\partial L_\tau}{\partial y_\tau} \phi_{out}^{\prime}(h_\tau) \right)^T
        \end{equation}
        \paragraph{反向传播}
        为了计算\ref{eq:1}，我们从末时刻($t = \tau$)向初时刻($t = 1$)计算，并且分为如下两步：
        \begin{enumerate}
            \item \textbf{$p_t$计算}：
                \begin{enumerate}
                    \item 当$t = \tau$时，根据线性方程\footnote{$\left(\ptl{^2 E_t}{h_t^2}\right)^{-1}$为对称矩阵，转置等于其本身}
                        \begin{equation}
                            \left(\ptl{^2 E_\tau}{h_\tau^2}\right) p_\tau = -v_\tau
                        \end{equation}
                        解出$p_\tau$
                    \item 当$t < \tau$时，根据线性方程：
                        \begin{equation}
                            \left(\ptl{^2 E_t}{h_t^2}\right) p_t = \left(\ptl{E_{t+1}}{h_{t+1} \partial h_t}\right)^T p_{t+1}
                        \end{equation}
                        解出$p_t$，而上式右边通过如下公式计算：
                        \begin{equation}
                            \left(\ptl{E_{t+1}}{h_{t+1} \partial h_t}\right)^T p_{t+1} = \left(\ptl{g_{t+1}}{h_t}\right)^T p_{t+1} = \ptl{}{h_t}(\tilde{p}_{t+1}^T g_{t+1})
                        \end{equation}
                        其中$\tilde{p}$表示计算过程中该参量视为常数，具体计算时先计算括号内的标量，再通过计算图算出关于$h_t$的导数。
                \end{enumerate}

            \item \textbf{参数梯度计算}:
                利用公式：
                \begin{equation}
                    \partial_t L_\tau := \left(\ptl{^2 E_t}{h_t \partial \theta}\right)^T p_t = \ptl{}{\theta}\left(\tilde{p}_t^T g_t\right)
                \end{equation}
                计算出每个时刻的参数梯度，上式计算顺序也是先求出括号内的标量，再通过计算图求导。
        \end{enumerate}

        最终梯度为；
        \begin{equation}
            \frac{\partial L_\tau}{\partial \theta} = \sum_{t=1}^{\tau} \partial_t L_\tau
        \end{equation}

\section{RNN网络(深度展开)}

    \subsection{理论推导}
    设RNN网络的迭代公式为：
    \begin{align}
        h_{t, 0} &= h_t\\
        h_{t, n+1} &= h_{t, n} - \alpha_{t, n} \ptl{E(h_{t, n}, h_t, x_t;\theta)}{h_{t, n}}\\
        h_{t+1} &= h_{t, K}
    \end{align}
    其中：
        \begin{itemize}
            \item $h_t$为t时刻的隐变量
            \item E为系统的能量函数
            \item $x_t$为系统t时刻的输入
            \item $\theta$为系统的可训练参数
            \item K为中间迭代步数上限
            \item $\alpha$为迭代步长控制参数
        \end{itemize}

    \subsection{算法概述}
        工程实现中，直接用pytorch的计算图构建功能，自动求得损失关于参数的导数。
        同时为了防止显存溢出，可能需要做以下工程处理：
        \begin{enumerate}
            \item \textbf{RNN传播截断：} 计算图中只保留RNN网络的最后若干步
            \item \textbf{梯度检查：}在前向传播计算那 $K$ 步下降时，不保存中间的激活值，而在反向传播需要计算梯度时，重新计算这 $K$ 步的前向过程来临时获取激活值。
        \end{enumerate}

\section{$RL^2$算法与PPO}
    

\end{document}